\documentclass[11pt]{article}
\usepackage[margin=25mm]{geometry}
\usepackage{fontspec}
\setmainfont{Latin Modern Roman}
\usepackage{amsmath,amssymb,mathtools}
\usepackage{unicode-math}
\setmathfont{Latin Modern Math}
\usepackage{xurl,hyperref}
\hypersetup{colorlinks=true,urlcolor=blue}
\setcounter{secnumdepth}{0}
\setlength{\parindent}{0pt}
\setlength{\parskip}{5pt}
\setlength{\emergencystretch}{3em}
\newcommand{\sref}[2]{\hyperlink{source:#1:#2}{[S#2]}}
\newcommand{\cataloguescope}[1]{\paragraph{Recorded target.}#1}
\newcommand{\R}{\mathbb R}
\newcommand{\C}{\mathbb C}
\newcommand{\N}{\mathbb N}
\newcommand{\Z}{\mathbb Z}
\newcommand{\Q}{\mathbb Q}
\newcommand{\F}{\mathbb F}
\newcommand{\eps}{\varepsilon}
\begin{document}
% UnsolvedMath ID 30002739; source catalogue code OWR-13355-003
\setcounter{section}{30002738}
\section{Periodic Gibbs Phases in Two Dimensions}\label{30002739}
{\small UnsolvedMath ID 30002739 · Mathematical Physics}
\subsection{Definitions and mathematical statement}

\paragraph{Shared notation.}$\mathbb N=\{1,2,\ldots\}$; the standard number fields have their usual meanings. All additional conventions are stated locally.


\paragraph{Definitions and conventions.}
For a lattice interaction of finite range, a Gibbs measure satisfies the Dobrushin--Lanford--Ruelle conditional-probability equations at inverse temperature $\beta\in(0,\infty)$. An extremal Gibbs measure is not a nontrivial convex mixture of other Gibbs measures for that same interaction. Periodicity means invariance under a finite-index subgroup of lattice translations, allowing a finite unit cell. Periodic is stronger than merely having a periodic-looking local correlation pattern.

\paragraph{Target retained from the canonical source catalogue.}
\paragraph{Catalogue wording retained for traceability.}
At positive temperature, is every extremal Gibbs measure for a finite-range interaction in two dimensions periodic?

\paragraph{Source scope and qualification.}
The allowed state space and interaction assumptions must be those of the report. In particular, the source statement should not be generalized without review to all continuous-spin, noncompact, or arbitrary inhomogeneous interactions.
\subsection{Short English statement}
Must every pure equilibrium phase of a finite-range two-dimensional system at positive temperature be periodic?
\subsection{Research attempt}
\paragraph{Model, date and investigated scope.}
OpenAI GPT-6 Astra (configured model \texttt{gpt-6-astra}), ultra reasoning; generated 11 October 2026. Proves the source property for finite-state on-site interactions and reduces the general uniqueness regime to translation covariance.

\paragraph{Sources and literature check.}
The complete supplied definition, including source qualifications and the dated literature assessment, was reviewed. The calculation below is an elementary restricted attempt motivated by [S2]; it does not claim novelty or a fresh literature-status determination.

\paragraph{Context and investigated scope.} The complete supplied definition, including source qualifications and the dated literature assessment, was reviewed. The calculation below is an elementary restricted attempt motivated by [S2]; it does not claim novelty or a fresh literature-status determination.

\paragraph{A completely verified periodic subcase.} Let spins take values in a finite set $S$ and let the translation-invariant interaction consist only of an on-site energy $h:S\to\mathbb R$. At inverse temperature $0<\beta<\infty$, put $p(s)=e^{-\beta h(s)}/\sum_{r\in S}e^{-\beta h(r)}$. The DLR conditional law on any finite region $\Lambda\subset\mathbb Z^2$ is $\prod_{x\in\Lambda}p(\sigma_x)$, independent of the exterior. Integrating this identity shows every Gibbs measure has these same finite-dimensional marginals. Therefore the Gibbs measure is uniquely the product measure and is translation invariant, hence periodic and extremal.

\paragraph{What uniqueness would imply generally.} For any translation-invariant specification, a lattice translate of a Gibbs measure is again Gibbs: translating each finite-region conditional equation gives the same specification on the translated region. If the Gibbs measure is unique, this forces invariance under every translation. Without uniqueness, extremality alone does not imply that its translate equals itself; an affine automorphism can permute different extreme points.

\paragraph{Why the hard question survives.} These proofs cover the decoupled model and the abstract uniqueness regime. They do not control interacting coexistence phases, where the source asks whether the translation orbit of every pure phase is finite. Inhomogeneous fields could trivially produce nonperiodic product laws, but they are outside the intended source assumptions and are not counterexamples. The first gap is a two-dimensional finite-range argument restricting translation orbits of extremal phases without assuming uniqueness or enlarging the state-space class.

\paragraph{Verification and adversarial audit.}
Checked exact DLR marginals, uniqueness/extremality and translated specifications; excluded inhomogeneous or unsupported state-space counterexamples. Model self-audit only; no independent human proof checking. Prompts A and B were applied to this canonical record and its qualified source scope. Checked exact DLR marginals, uniqueness/extremality and translated specifications; excluded inhomogeneous or unsupported state-space counterexamples. Retained unresolved status and zero original-target progress; No bound on translation orbits of extremal interacting coexistence phases is proved.

Coordinator checked the on-site DLR uniqueness proof and the implication from uniqueness to full translation invariance. This additional check was model-assisted, not independent human verification.

\paragraph{Tools, process and contributions.}
Python 3 was used for file preparation and structural checks; no numerical experiment or theorem prover was used. Contribution prompt A was followed by prompt B. GPT-6 Astra with ultra reasoning drafted the research and performed its stated self-audit; any additional coordinator contribution is identified in the verification and audit. A GPT-6-based Codex coordinator assembled the manuscript and checked repository integration; its exact variant and reasoning configuration were not exposed. No human mathematical contribution, independent human proof check or formal proof certification is supplied. Complete derivations and the consulted source references are included here; private reasoning transcripts are not published.

\paragraph{Final status and first remaining gap.}
\textbf{UNRESOLVED}. No bound on translation orbits of extremal interacting coexistence phases is proved.

\section{Completion Estimate}
\noindent\textbf{COMPLETION: 0\%}
The calculation is an auxiliary or restricted attempt and establishes no newly verified portion of the full source target. No bound on translation orbits of extremal interacting coexistence phases is proved.

\subsection{Sources}
\begin{itemize}
\item[\hypertarget{source:30002739:1}{S1}] ULAM.AI and the UnsolvedMath Contributors, UnsolvedMath, record 30002739 (OWR-13355-003). Catalogue curation under CC BY 4.0; definitions and English exposition added locally.
\url{https://huggingface.co/datasets/ulamai/UnsolvedMath}.
\item[\hypertarget{source:30002739:2}{S2}] Mini-Workshop: Dynamical versus Diffraction Spectra in the Theory of Quasicrystals (2014), Oberwolfach Reports, DOI 10.4171/owr/2014/53, pp. 2987--3013. Primary report accessed 2026-10-11; relevant PDF pages 9, 10.
\url{https://doi.org/10.4171/owr/2014/53}.
\item[\hypertarget{source:30002739:3}{S3}] Koch and Radin, positive-temperature quasicrystal model (catalogue-listed context; not independently consulted in this preparation).
\url{https://doi.org/10.1007/s10955-009-9896-9}.
\end{itemize}
\label{end:30002739}
\end{document}
